\title{DeepSeekMath: Pushing the Limits of Mathematical Reasoning in Open Language Models}

\begin{abstract}
Mathematical reasoning remains a considerable obstacle for language models, given the inherent intricacies and structured demands of the domain. Here, we present DeepSeekMath~7B, which is derived by continuing the pre-training of DeepSeek-Coder-Base-v1.5 7B on a corpus of 120B math-centric tokens extracted from Common Crawl, supplemented by natural language and programming data. DeepSeekMath~7B attains a notable 51.7\% on the MATH benchmark—a measure designed for advanced mathematical problem-solving—without utilizing supplementary toolkits or ensemble methods, narrowing the gap with closed models such as Gemini-Ultra and GPT-4. Utilizing self-consistency sampling (with 64 samples), DeepSeekMath~7B reaches 60.9\% on MATH. This enhancement in mathematical reasoning stems from two principal innovations: first, a carefully designed data selection pipeline that leverages extensive open web data; second, the introduction of Group Relative Policy Optimization (GRPO), an adaptation of Proximal Policy Optimization (PPO), which both elevates mathematical reasoning proficiency and improves the efficiency of PPO's memory usage.
\end{abstract}

\section{Introduction}

Large language models (LLMs) have transformed approaches to mathematical reasoning within artificial intelligence, leading to notable progress across quantitative reasoning \citep{MATH} and geometry reasoning \citep{trinh2024solving} benchmarks. Additionally, these systems have served as valuable resources for human users engaged in intricate mathematical problem-solving tasks \citep{tao}. Despite these advances, top-performing solutions such as GPT-4 \citep{gpt4} and Gemini-Ultra \citep{gemini} remain closed-source, and the current generation of open-source alternatives significantly lags in terms of their capabilities.

This paper presents DeepSeekMath, a specialized language model tailored for mathematical tasks, which not only surpasses existing open-source models in mathematical reasoning but also closely matches the performance of GPT-4 on established academic benchmarks. Central to this achievement is the development of the DeepSeekMath Corpus, an extensive, high-quality pre-training dataset containing 120B math tokens. The data is sourced from Common Crawl (CC), employing a classifier based on fastText \citep{joulin2016fasttext}. In its initial phase, this classifier is trained with positive samples from OpenWebMath \citep{openwebmath} and a varied set of negative samples drawn from general web content. The classifier is then used to identify additional relevant samples from CC, which are subsequently filtered and validated by human annotators. These newly annotated examples are integrated to further retrain the classifier, thereby enhancing its accuracy. Performance evaluations demonstrate that this comprehensive dataset leads to strong model performance: DeepSeekMath-Base 7B attains 64.2\% accuracy on GSM8K \citep{gsm8k} and 36.2\% on the MATH competition benchmark \citep{MATH}, outperforming models such as Minerva 540B \citep{minerva}. Furthermore, because the DeepSeekMath Corpus is multilingual, notable gains are observed on Chinese mathematical evaluations \citep{wei2023cmath,agieval}. The methodologies and insights from our approach to mathematical data curation are intended as a foundation for ongoing and future research, as considerable opportunities remain for further advancements.

DeepSeekMath-Base is initialized using DeepSeek-Coder-Base-v1.5 7B \citep{deepseek-coder}, since we found that leveraging a model pre-trained on code yields superior results relative to initiating from a standard large language model. Additionally, we observed that math-focused training leads to improvements not only in the model’s mathematical reasoning, but also enhances general reasoning skills, as evidenced by higher scores on the MMLU \citep{mmlu} and BBH \citep{bbh} benchmarks.

Following pre-training, DeepSeekMath-Base undergoes instruction tuning for mathematics tasks, drawing on chain-of-thought \citep{cot}, program-of-thought \citep{pot,pal}, and tool-integrated reasoning \citep{tora} datasets. The resulting DeepSeekMath-Instruct 7B demonstrates superior performance compared to other 7B-sized models, and achieves results on par with open-source instruction-tuned models in the 70B parameter range.

In addition, we present Group Relative Policy Optimization (GRPO), an alternative reinforcement learning (RL) approach related to Proximal Policy Optimization (PPO) \citep{schulman2017proximal}. Unlike standard PPO, GRPO eliminates the need for a critic model by utilizing group-based score baselines, which leads to significant savings in training resources. By applying GRPO to only a subset of English instruction tuning data, we observe marked gains over the strong baseline of DeepSeekMath-Instruct in both in-domain (GSM8K: 82.9\% $\rightarrow$ 88.2\%, MATH: 46.8\% $\rightarrow$ 51.7\%) and out-of-domain (e.g., CMATH: 84.6\% $\rightarrow$ 88.8\%) math tasks during RL training. 

Moreover, we offer a comprehensive framework for interpreting methods such as Rejection Sampling Fine-Tuning (RFT) \citep{yuan2023scaling}, Direct Preference Optimization (DPO) \citep{dpo}, PPO, and GRPO, revealing that these strategies can be viewed as direct or streamlined forms of RL. Through a range of systematic experiments—including comparisons of online and offline learning, analysis of outcome versus process supervision, and evaluations of single-turn versus iterative RL—we dissect the fundamental factors that underlie this unified framework. Finally, we elucidate the mechanisms by which our RL method enhances the performance of instruction-tuned models and outline prospective directions for achieving even more efficient RL based on this integrative approach.

\subsection{Contributions}
This work advances the field by introducing large-scale math pre-training techniques and conducting a detailed investigation into reinforcement learning.

\textbf{Math Pre-Training at Scale}

\begin{itemize}[topsep=0pt]
    \item We demonstrate that publicly available Common Crawl data can be effectively harnessed for mathematical tasks. Through the development of a rigorous data filtering and selection framework, we assemble the DeepSeekMath Corpus—a high-quality collection containing 120B tokens, sourced from mathematically relevant web content. This dataset substantially surpasses prior efforts, being nearly seven times larger than that used by Minerva \citep{minerva} and over nine times greater than the dataset of OpenWebMath \citep{openwebmath}.

    \item Our model, DeepSeekMath-Base 7B, attains performance on par with Minerva 540B \citep{minerva}, underscoring that model scale, in terms of parameters, is not the sole determinant of mathematical reasoning proficiency. High performance can also be realized by smaller models pre-trained on robust datasets.

    \item We provide insights gleaned from our mathematical pre-training experiments, notably that preceding math training with code-focused training enhances models’ aptitude for solving mathematical problems—regardless of whether tools are used. This observation partially addresses the ongoing debate: \textit{does code training contribute to better reasoning skills?} Our findings support this for mathematical reasoning in particular.

    \item While incorporating arXiv papers into pre-training is a widespread practice, especially among mathematics-focused studies, our results indicate that doing so does not yield meaningful gains across the mathematical benchmarks evaluated in this study.
\end{itemize}

\textbf{Exploration and Analysis of Reinforcement Learning}

\begin{itemize}[topsep=0pt]
    \item We present Group Relative Policy Optimization (GRPO), a resource-efficient reinforcement learning framework. Rather than relying on a critic network, GRPO utilizes group-derived scores to estimate the baseline, markedly lowering computational demands in contrast to Proximal Policy Optimization (PPO).
    \item Our results indicate that GRPO substantially boosts the effectiveness of our instruction-tuned system, DeepSeekMath-Instruct, using only instruction-tuning datasets. Additionally, we detect gains in generalization to out-of-domain tasks throughout the reinforcement learning process.
    \item We propose a comprehensive framework for conceptualizing various approaches—including RFT, DPO, PPO, and GRPO. Through a series of systematic experiments (including comparisons between online and offline training, outcome-based versus process-based supervision, and single-turn versus iterative reinforcement learning), we dissect the critical features underpinning this framework.
    \item Utilizing the proposed unified paradigm, we investigate the mechanisms driving reinforcement learning success and outline multiple promising avenues for advancing the efficiency of LLM reinforcement learning.
\end{itemize}

\subsection{Summary of Evaluations and Metrics}

\begin{itemize}[topsep=0pt]
    \item \textbf{English and Chinese Mathematical Reasoning}: Our evaluation encompasses an extensive suite of English and Chinese datasets, spanning mathematical reasoning tasks from elementary through collegiate levels. For English, the benchmarks considered are GSM8K \citep{gsm8k}, MATH \citep{MATH}, SAT \citep{llemma}, OCW Courses \citep{minerva}, and MMLU-STEM \citep{mmlu}. Chinese benchmarks include MGSM-zh \citep{mgsm}, CMATH \citep{wei2023cmath}, Gaokao-MathCloze \citep{agieval}, and Gaokao-MathQA \citep{agieval}. We assess the models' capability to independently generate detailed text solutions without external tools, alongside their problem-solving aptitude with Python assistance.

    DeepSeekMath-Base rivals the proprietary Minerva 540B model \citep{minerva} on English tasks and exceeds the performance of all open-source counterparts (such as Mistral 7B \citep{mistral} and Llemma-34B \citep{llemma}), regardless of their prior exposure to math pre-training, frequently achieving significantly better results. Particularly in Chinese tasks, DeepSeekMath-Base delivers superior outcomes, a result we attribute to incorporating both English and high-quality non-English mathematical data, contrary to previous work \citep{minerva,llemma} that limited pre-training to English sources. With subsequent application of instruction tuning and reinforcement learning, the developed models—DeepSeekMath-Instruct and DeepSeekMath-RL—reach new benchmarks, becoming the first open-source systems to surpass 50\% accuracy on the highly challenging MATH dataset.
\end{itemize}

\item \textbf{Formal Mathematics}: We assess DeepSeekMath-Base on the informal-to-formal theorem proving benchmark from \citep{dsp_proof}, specifically utilizing miniF2F \citep{minif2f} with Isabelle \citep{isabelle} serving as the proof assistant. DeepSeekMath-Base exhibits impressive few-shot capabilities for autoformalization.

\item \textbf{Natural Language Understanding, Reasoning, and Code}: To obtain a broad assessment of the models' proficiency in language comprehension, logical reasoning, and programming, we test DeepSeekMath-Base on several benchmarks: Massive Multitask Language Understanding (MMLU) \citep{mmlu}, which features 57 multiple-choice topics across various domains; BIG-Bench Hard (BBH) \citep{bbh}, a set of 23 rigorously multi-step reasoning tasks; and the code-centric benchmarks HumanEval \citep{codex} and MBPP \citep{mbpp}, both commonly used for evaluating code language models. Pre-training on mathematics proves advantageous for both comprehension and reasoning scores.

\section{Math Pre-Training}

\subsection{Data Collection and Decontamination} This section describes the methodology for building the DeepSeekMath Corpus sourced from Common Crawl. As illustrated in Figure \ref{fig:data_collect}, an iterative process is employed to amass a comprehensive mathematical dataset beginning with an initial seed (for example, a select high-quality set of mathematical data). This methodology is also transferable to other areas, such as programming languages.

The process begins by adopting OpenWebMath \citep{openwebmath}—a vetted collection of mathematical text from the web—as the foundational seed corpus. Utilizing this dataset, we train a fastText model \citep{joulin2016fasttext} to identify and extract additional web pages similar to those in OpenWebMath. For model training, 500,000 randomly chosen items from the seed corpus serve as positive samples, and another 500,000 web pages sourced from Common Crawl serve as negative examples. Model training is conducted using an open-source package\footnote{\url{https://fasttext.cc}}, with specific settings: vector dimensionality of 256, learning rate set to 0.1, a word n-gram maximum length of 3, a minimum word frequency of 3, and 3 epochs of training. To manage the dataset size from Common Crawl, both URL-based and near-duplicate removal procedures are applied, narrowing it to 40 billion HTML documents. The trained fastText model is then leveraged to extract math-relevant pages from the deduplicated dataset. Subsequently, all retrieved pages are ranked according to the fastText model's predicted relevance, retaining only the highest-scoring items. The preserved dataset size is determined through pre-training trials using the top 40B, 80B, 120B, and 160B tokens. For the initial cycle, the upper 40B tokens are retained.

Following the initial round of data collection, a significant portion of mathematical web pages remains uncaptured. This shortfall is primarily attributed to the insufficient diversity within the set of positive samples used to train the fastText model. To address this, we incorporate more mathematical web resources to expand the seed corpus, thereby enhancing the model’s ability to identify mathematical content. Specifically, we segment the entire Common Crawl dataset into non-overlapping domains, each defined as a set of web pages sharing a common base URL. For each domain, we determine the proportion of pages that were collected in the first round. Domains with more than 10\% of their pages included are designated as math-related (for instance, \url{mathoverflow.net}).

We then conduct a manual review to label URLs within these domains that point to mathematical content (such as \url{mathoverflow.net/questions}). Any uncollected web pages linked to these URLs are subsequently incorporated into the seed corpus. This expanded pool of positive examples allows us to train an improved fastText model, increasing its recall of mathematical data in subsequent rounds. After four collection cycles, we have amassed a total of 35.5M mathematical web pages, comprising 120B tokens. During the fourth iteration, we observe that almost 98\% of the data was already acquired in the previous cycle, leading us to halt further collection efforts.

To prevent contamination from benchmark datasets, we apply the filtering method used in \cite{deepseek-coder}, excluding any web pages that contain questions or answers sourced from major mathematical benchmarks. This includes English datasets like GSM8K \citep{gsm8k} and MATH \citep{MATH}, as well as Chinese benchmarks such as CMATH \citep{wei2023cmath} and AGIEval \citep{agieval}. Our filtering rule discards any text that contains an exact match to a 10-gram substring found in these benchmark sets. For benchmark passages shorter than 10 grams but containing at least 3 grams, we use exact string matching to eliminate contaminated content from the math training data.

\subsection{Validating the Quality of the DeepSeekMath~Corpus}

To assess the value of the DeepSeekMath~Corpus, we conduct pre-training experiments comparing it with recently published math-oriented corpora:

\begin{itemize}[topsep=0pt]
    \item \textbf{MathPile} \citep{mathpile}: A composite dataset (8.9B tokens) compiled from sources such as textbooks, Wikipedia, ProofWiki, CommonCrawl, StackExchange, and arXiv, with the bulk (over 85\%) drawn from arXiv;
    \item \textbf{OpenWebMath} \citep{openwebmath}: Extracted from CommonCrawl by filtering for mathematical materials, containing 13.6B tokens;
    \item \textbf{Proof-Pile-2} \citep{llemma}: Encompassing OpenWebMath, AlgebraicStack (10.3B tokens focused on mathematical code), and arXiv papers (28.0B tokens). In our experiments on Proof-Pile-2, we adhere to the arXiv:Web:Code mixture ratio of 2:4:1, as specified in \cite{llemma}.
\end{itemize}

\subsubsection{Training Setting}
\label{sec:quality-policy}
We conduct mathematical training on a broadly pre-trained language model comprising 1.3 billion parameters, employing the same architecture as DeepSeek LLMs \citep{deepseek-llm}, henceforth referred to as DeepSeek-LLM 1.3B. Separate training is carried out on each mathematical dataset for a total of 150 billion tokens. The experiments utilize the resource-efficient HAI-LLM \citep{haillm} training infrastructure. In alignment with DeepSeek LLM training protocols, the AdamW optimizer \citep{adamW} is employed, configured with $\beta_1=0.9$, $\beta_2=0.95$, and $\mathrm{weight\_decay}=0.1$. A multi-stage learning rate schedule is applied: the learning rate is increased to its maximum following 2,000 steps of linear warmup, reduced to 31.6\% of its peak at 80\% completion of training, and further decreased to 10.0\% of the initial peak at 90\% of the training duration. The peak learning rate is set at 5.3e-4. Training employs a token batch size of 4 million, with a maximum context window of 4K tokens.

\subsubsection{Evaluation Results}

\textbf{The DeepSeekMath~Corpus stands out for its exceptional quality, extensive multilingual mathematical data, and unprecedented scale.}

\begin{itemize}[topsep=0pt]
    \item \textbf{High-quality}:
    We assess the impact on downstream tasks across 8 mathematical benchmarks using few-shot chain-of-thought prompting \cite{cot}. As detailed in Table \ref{tab:corpora_comparison}, the model trained with the DeepSeekMath~Corpus consistently surpasses others in performance. Further, Figure \ref{fig:corpora_comparisons} illustrates that a model trained with the DeepSeekMath Corpus exceeds the performance of Proof-Pile-2 even after only 50B tokens (equivalent to one complete epoch over Proof-Pile-2), suggesting that the DeepSeekMath Corpus exhibits superior average quality.
    \item \textbf{Multilingual}:
    The DeepSeekMath~Corpus integrates resources across several languages, with English and Chinese constituting the bulk of its content. Table \ref{tab:corpora_comparison} demonstrates that training on the DeepSeekMath~Corpus yields enhanced mathematical reasoning abilities in both English and Chinese, whereas models relying on other largely English-centric corpora show little improvement or may even degrade Chinese reasoning capabilities.
    \item \textbf{Large-scale}:
    Compared to currently available mathematical corpora, the DeepSeekMath~Corpus is substantially larger. Figure \ref{fig:corpora_comparisons} indicates that DeepSeek-LLM 1.3B, trained on this dataset, attains a sharper learning trajectory and more sustained advancements. By contrast, alternative corpora, being much smaller, are cycled through multiple epochs, causing the respective model's progress to stagnate quickly.
\end{itemize}

\subsection{Training and Evaluating DeepSeekMath-Base 7B}

This section details DeepSeekMath-Base 7B, a foundational model characterized by robust reasoning capacities, particularly within mathematical domains. The model builds upon the DeepSeek-Coder-Base-v1.5 7B architecture \citep{deepseek-coder}, using it as the initialization point, and is subsequently trained on a total of 500B tokens. The dataset composition comprises 56\% DeepSeekMath Corpus, 4\% AlgebraicStack, 10\% arXiv, 20\% Github code, and 10\% natural language data from Common Crawl, sourced in both English and Chinese. With respect to training configuration, we generally adhere to the methodology outlined in Section \ref{sec:quality-policy}, with the exception of setting the learning rate upper limit to 4.2e-4 and employing a batch size corresponding to 10M tokens.

To thoroughly examine the mathematical proficiency of DeepSeekMath-Base 7B, we assess the model's capacity for generating comprehensive mathematical solutions without external assistance, its effectiveness in addressing mathematical problems utilizing auxiliary tools, and its aptitude for formal theorem verification. In addition to mathematical tasks, we also survey the base model's overall capabilities, specifically regarding natural language understanding, logical reasoning, and programming aptitude.

\paragraph{Mathematical Problem Solving with Step-by-Step Reasoning}

DeepSeekMath-Base's aptitude for solving mathematical questions is tested using few-shot chain-of-thought prompting \citep{cot}, spanning eight evaluation benchmarks in both English and Chinese. These benchmarks target quantitative reasoning (for instance, GSM8K \citep{gsm8k}, MATH \citep{MATH}, and CMATH \citep{wei2023cmath}) as well as multiple-choice tasks (such as MMLU-STEM \citep{mmlu} and Gaokao-MathQA \citep{agieval}), thereby encompassing a broad spectrum of mathematical disciplines from elementary through to undergraduate complexity.

As illustrated in Table \ref{tab:base_math_cot}, among all open-source base models, DeepSeekMath-Base 7B achieves top results across the eight evaluated benchmarks. This includes outperforming prominent models such as the widely-adopted generalist Mistral 7B \citep{mistral} and the more recent Llemma 34B \citep{llemma}, which was trained on Proof-Pile-2 \citep{llemma} for mathematics tasks. Remarkably, DeepSeekMath-Base not only surpasses all existing open-source models by more than 10\% absolute accuracy on the advanced MATH dataset, but also outdoes Minerva 540B \citep{minerva}, a closed-source model of considerably larger scale (77 times the parameters), which extends PaLM \citep{palm} through further mathematical training.

\paragraph{Mathematical Problem Solving with Tool Use} For evaluating program-supported mathematical reasoning, we test on GSM8K and MATH using a few-shot program-of-thought prompting framework \citep{pot,pal}. In this setting, the model is guided to construct a Python script, leveraging libraries such as \textit{math} and \textit{sympy}, to handle complex calculations for each given problem. The output is then determined by running the generated program. Table \ref{tab:base_math_pot_proof} presents results where DeepSeekMath-Base 7B exceeds the previously leading Llemma 34B in performance.

\paragraph{Formal Mathematics}

Automation of formal proof writing serves to increase precision and reliability in mathematical arguments, contributing to greater efficiency—a topic attracting increasing research attention. We assess DeepSeekMath-Base 7B using the informal-to-formal proof translation task described in \citep{dsp_proof}: the aim is to produce a formal proof from an informal problem statement, its formalized equivalent, and a corresponding informal proof. Our evaluation leverages miniF2F \citep{minif2f}, a benchmark focused on formal mathematics at Olympiad-level difficulty, with models generating Isabelle formal proofs via few-shot prompting. In line with \cite{dsp_proof}, proof sketches are first created by the model, after which Sledgehammer \citep{sledgehammer} is used to automatically complete any missing proof steps. Table \ref{tab:base_math_pot_proof} demonstrates that DeepSeekMath-Base 7B is highly effective at autoformalization of proofs.

\paragraph{Natural Language Understanding, Reasoning, and Code}

Model ability on language comprehension is measured by MMLU \citep{mmlu}, while reasoning is assessed with BBH \citep{bbh}, and code generation on HumanEval \citep{codex} and MBPP \citep{mbpp}. Results shown in Table \ref{tab:understanding_reasoning_code} indicate that DeepSeekMath-Base 7B achieves substantial improvements on MMLU and BBH over its predecessor DeepSeek-Coder-Base-v1.5 \citep{deepseek-coder}, underscoring the advantages conferred by mathematical training for language understanding and reasoning. Furthermore, through the incorporation of code token training, DeepSeekMath-Base 7B is able to preserve the coding proficiency of DeepSeek-Coder-Base-v1.5 on the code-related benchmarks. Across the reasoning and code generation tasks, DeepSeekMath-Base 7B consistently surpasses the general-purpose Mistral 7B \citep{mistral}.

\section{Supervised Fine-Tuning}

\subsection{SFT Data Curation}

We assemble a dataset for instruction-based mathematical fine-tuning that spans both English and Chinese questions drawn from a variety of mathematical domains and difficulty levels. Each question is aligned with solutions employing chain-of-thought (CoT) \citep{cot}, program-of-thought (PoT) \citep{pot,pal}, or tool-integrated reasoning approaches \citep{tora}. Our final training set contains 776K examples.

\begin{itemize}[topsep=0pt]
    \item \textbf{English mathematical datasets}: We enrich GSM8K and MATH with annotated tool-assisted answers, and further incorporate a subset of MathInstruct \citep{MathInstruct} and the Lila-OOD training split \citep{lila}, utilizing either CoT or PoT methodologies for problem solving. This English dataset represents a broad array of mathematical disciplines, such as algebra, probability, number theory, calculus, and geometry.
    \item \textbf{Chinese mathematical datasets}: We gather a collection of K-12 mathematics problems in Chinese covering 76 distinct subfields, including linear equations, and annotate solutions using both CoT and tool-assisted reasoning.
\end{itemize}

\subsection{Training and Evaluating DeepSeekMath-Instruct 7B}

Here we present the procedure for training DeepSeekMath-Instruct 7B, which is instruction-tuned for mathematics on top of DeepSeekMath-Base. To generate training batches, multiple samples are concatenated until reaching the 4K token context window limit. The model is optimized over 500 steps, utilizing a batch size of 256 and a fixed learning rate of 5e-5.

To measure the model's capabilities in mathematics, we conduct assessments on four quantitative reasoning benchmarks, both in English and Chinese, considering settings with and without tool usage. We compare our system's outcomes with those of contemporaneous leading models:

\begin{itemize}[topsep=0pt]
    \item \textbf{Closed-source models} evaluated include: (1) members of the GPT series, with GPT-4 \citep{gpt4} and GPT-4 Code Interpreter~\footnote{\url{https://openai.com/blog/chatgpt-plugins\#\#code-interpreter}} being the most advanced representatives, (2) Gemini Ultra and Pro \citep{gemini}, (3) Inflection-2 \citep{inflection-2}, (4) Grok-1~\footnote{\url{https://x.ai/model-card}},
    as well as recent releases from Chinese industry, such as (5) Baichuan-3~\footnote{\url{https://www.baichuan-ai.com}},
    and (6) the most recent member of the GLM series, GLM-4~\footnote{\url{https://open.bigmodel.cn/dev/api\#glm-4}}

    from the GLM family \citep{glm}.
\end{itemize}

The models discussed serve broad applications and have generally been subjected to comprehensive alignment processes.

\item \textbf{Open-source models} encompass the following: foundational models such as (1) DeepSeek-LLM-Chat 67B \citep{deepseek-llm}, (2) Qwen 72B \citep{qwen}, (3) SeaLLM-v2 7B \citep{seallm}, and (4) ChatGLM3 6B \citep{chatglm3}, alongside math-augmented variants: (5) InternLM2-Math 20B\footnote{\url{https://github.com/InternLM/InternLM-Math}}, based on InternLM2 and subjected to mathematical pre-training with subsequent instruction tuning, (6) Math-Shepherd-Mistral 7B, which employs PPO training \citep{schulman2017proximal} on Mistral 7B \citep{mistral} utilizing a process-supervised reward model, (7) the WizardMath suite \citep{wizardmath}, which enhances mathematical problem-solving for both Mistral 7B and Llama-2 70B \citep{llama2} through evolve-instruct (a variant of instruction tuning reliant on instructions generated by AI) and PPO, trained primarily with problems from GSM8K and MATH, (8) MetaMath 70B \citep{metamath}, which consists of Llama-2 70B further fine-tuned with an enriched version of GSM8K and MATH, (9) ToRA 34B \cite{tora}, representing CodeLlama 34B specialized via fine-tuning for mathematical reasoning augmented by tool usage, and (10) MAmmoTH 70B \citep{MathInstruct}, which builds on Llama-2 70B further tuned on MathInstruct.

Table \ref{tab:sft_rl_math} summarizes model performance when tool assistance is restricted. DeepSeekMath-Instruct 7B achieves robust results in step-wise mathematical reasoning. Remarkably, on the competition-grade MATH benchmark, this model exceeds all open-source alternatives and most proprietary models (e.g., Inflection-2 and Gemini Pro) by an absolute margin of at least 9\%, including models of significantly greater size (e.g., Qwen 72B) or those that underwent math-centric RL finetuning (such as WizardMath-v1.1 7B). Although DeepSeekMath-Instruct's performance is on par with leading Chinese proprietary models (GLM-4 and Baichuan-3) on MATH, it still trails behind top-tier commercial models like GPT-4 and Gemini Ultra.

In scenarios permitting both natural language reasoning and the use of external tools via programming to solve problems, DeepSeekMath-Instruct 7B attains an accuracy close to 60\% on the MATH benchmark—outperforming all publicly available models. Across additional benchmarks, it maintains competitiveness with DeepSeek-LLM-Chat 67B, which, at tenfold the size, previously held the state-of-the-art.

\section{Reinforcement Learning}

\subsection{Group Relative Policy Optimization} Reinforcement learning (RL) has demonstrated efficacy in enhancing the mathematical reasoning skills of large language models beyond the Supervised Fine-Tuning (SFT) phase \citep{wang2023math,wizardmath}. This section introduces our RL approach, Group Relative Policy Optimization (GRPO), which is both efficient and effective.

\subsubsection{From PPO to GRPO} Proximal Policy Optimization (PPO) \citep{schulman2017proximal}, a widely-adopted actor-critic RL algorithm, is frequently utilized for the RL fine-tuning of LLMs \citep{ouyang2022training}. Specifically, PPO seeks to optimize the language model by maximizing this surrogate objective function:

\begin{equation}
\footnotesize
    \mathcal{J}_{PPO}(\theta) = \mathbb{E}{[q \sim P(Q), o \sim \pi_{\theta_{old}}(O|q)]} \frac{1}{|o|} \sum_{t=1}^{|o|} \min \left[ \frac{\pi_\theta(o_{t} | q, o_{<t})}{\pi_{\theta_{old}}(o_{t} | q, o_{<t})} A_{t}, \text{clip} \left( \frac{\pi_\theta(o_{t} | q, o_{<t})}{\pi_{\theta_{old}}(o_{t} | q, o_{<t})}, 1 - \varepsilon, 1

Here, $\pi_{\theta}$ and $\pi_{\theta_{old}}$ denote the current and previous iterations of the policy models, respectively, while $q$ and $o$ correspond to questions and outputs, sampled from both the question set and from $\pi_{\theta_{old}}$. The scalar $\varepsilon$ is a PPO-specific hyper-parameter that controls clipping, thereby promoting stability in training. The advantage $A_t$ is derived using Generalized Advantage Estimation (GAE) \citep{gae}, which leverages the observed rewards $\{r_{\ge t}\}$ and an estimated value function $V_{\psi}$. Consequently, PPO requires the concurrent training of a value function model along with the policy model. To prevent excessive optimization toward the reward model, a conventional strategy involves incorporating a per-token KL divergence penalty between the output probabilities of the current policy and a reference model, applied at each token in the reward formulation \citep{ouyang2022training}:

\begin{equation}
    r_{t} = r_\varphi(q, o_{\le t}) - \beta \log\frac{\pi_{\theta}(o_{t}|q, o_{<t})}{\pi_{ref}(o_{t}|q, o_{<t})},
\label{eq:PPO-reward}
\end{equation}

where $r_\varphi$ designates the reward model, $\pi_{ref}$ refers to the reference model (commonly, the initial SFT model), and $\beta$ modulates the strength of the KL penalty.

Training a value function in PPO typically necessitates a model of similar scale to the policy, imposing significant resource demands in terms of memory and computation. In reinforcement learning settings, the value function primarily serves as a baseline to reduce variance in advantage estimates. However, for large language models, the reward model generally provides feedback only at the final token, complicating the design of a per-token accurate value function. To address these limitations, we introduce Group Relative Policy Optimization (GRPO) as illustrated in Figure \ref{fig:grpo}, which eliminates the need for a separate value function as required in PPO. Instead, GRPO employs the mean reward over several sampled outputs for the same question to establish the baseline. Specifically, for each input $q$, we sample a set of outputs $\{o_1, o_2, \cdots, o_G\}$ from $\pi_{\theta_{old}}$, and update the policy by maximizing the objective:

\begin{equation}
\footnotesize
\begin{split}
    \mathcal{J}_{GRPO}(\theta) &= \mathbb{E}{[q \sim P(Q), \{o_i\}_{i=1}^G \sim \pi_{\theta_{old}}(O|q)]}  \\
    & \frac{1}{G}\sum_{i=1}^G\frac{1}{|o_i|} \sum_{t=1}^{|o_i|} \left\{ \min \left[ \frac{\pi_\theta(o_{i,t} | q, o_{i,<t})}{\pi_{\theta_{old}}(o_{i,t} | q, o_{i,<t})} \hat{A}_{i,t}, \text{clip} \left( \frac{\pi_\theta(o_{i,t} | q, o_{i,<t})}{\pi_{\theta_{old}}(o_{i,t} | q, o_{i,<t})}, 1 - \varepsilon, 1 + \varepsilon \right)  \hat{A}_{i,t} \right] - \beta \mathbb{D}_{KL}\left[\pi_{\theta} || \pi_{ref}\right]\right\} ,
\end{split}
\label{eq:GRPO-obj}
\end{equation}

In this formulation, $\varepsilon$ and $\beta$ remain as tunable hyper-parameters, and $\hat{A}_{i,t}$ signifies the advantage, determined using only the relative rewards among the outputs in the same group—a process described in further detail below. GRPO's group-based approach to computing advantages is particularly congruent with the comparative format in which reward models are usually trained, as these models operate over collections of outputs related to the same question. Notably, GRPO differs from PPO in its regularization approach: instead of imposing a KL penalty on the reward signal, GRPO directly penalizes the KL divergence between the policy under training and the reference policy at the loss level, thus sidestepping added complexity in calculating $\hat{A}_{i,t}$. Also, contrary to the KL term found in (\ref{eq:PPO-reward}), we use an unbiased

\begin{algorithm}[t]
  \small
  \caption{Iterative Group Relative Policy Optimization}
  \textbf{Input} initial policy model $\pi_{\theta_{\text{init}}}$; reward models $r_\varphi$; task prompts $\mathcal{D}$; 
  hyperparameters $\varepsilon$, $\beta$, $\mu$
  \begin{algorithmic}[1]
    \State Set policy model $\pi_\theta \leftarrow \pi_{\theta_{\text{init}}}$
    \For{iteration = 1, \dots, I}
       \State Set reference model $\pi_{ref} \leftarrow \pi_{\theta}$
      \For{step = 1, \dots, M}
      \State Draw a minibatch $\mathcal{D}_b$ from $\mathcal{D}$
      \State Assign current policy to old policy: $\pi_{\theta_{old}} \leftarrow \pi_{\theta}$ 
      \State For each question $q \in \mathcal{D}_b$, generate $G$ outputs $\{o_i\}_{i=1}^G \sim \pi_{\theta_{old}} (\cdot \mid q) $
      \State For every sampled $o_i$, use the reward model $r_\varphi$ to obtain rewards $\{r_i\}_{i=1}^{G}$
      \State Estimate group-relative advantages $\hat{A}_{i,t}$ for the $t$-th token in each $o_i$
      \For{GRPO iteration = 1, \dots, $\mu$}
        \State Adjust $\pi_\theta$ to maximize the GRPO objective (Equation \ref{eq:GC-GRPO})
      \EndFor
    \EndFor 
    \State Continuously refine $r_\varphi$ utilizing a replay strategy.
    \EndFor 
  \end{algorithmic}
  \textbf{Output} $\pi_\theta$
  \label{alg:iter-grpo}
\end{algorithm}

\subsubsection{Outcome Supervision RL with GRPO} In a formal sense, for each given question $q$, a set of $G$ responses $\{o_1, o_2, \cdots, o_G\}$ is sampled from the previous version of the policy model $\pi_{\theta_{old}}$. The reward model assigns a scalar reward to every output, resulting in a vector of rewards $\mathbf{r}=\{r_1, r_2, \cdots, r_G\}$. Next, these reward values are standardized by centering and scaling with the group mean and standard deviation. Under outcome supervision, the normalized reward, calculated at the end of every output $o_i$, is used to assign an identical advantage $\hat{A}_{i, t}$ for every token in $o_i$; specifically, $\hat{A}_{i, t} = \widetilde{r}_i = \frac{r_i- {\rm mean}(\mathbf{r})}{{\rm std}(\mathbf{r})}$. The policy update is then performed by maximizing the objective in equation (\ref{eq:GRPO-obj}).

\subsubsection{Process Supervision RL with GRPO} Because outcome supervision restricts feedback to only terminal rewards, its ability to efficiently guide policy learning in sophisticated mathematical domains may be limited. To address this, inspired by \cite{wang2023math}, we incorporate process supervision, where rewards are delivered at each intermediate reasoning stage. Concretely, for each prompt $q$ and the batch of $G$ sampled responses $\{o_1, o_2, \cdots, o_G\}$, a process reward model evaluates each intermediate reasoning step, producing rewards organized as $\mathbf{R} = \{ \{r_1^{index(1)},\cdots,r_1^{index(K_1)}\}, \cdots,  \{r_G^{index(1)},\cdots,r_G^{index(K_G)}\} \}$, with $index(j)$ denoting the terminal token index for the $j$-th step and $K_i$ specifying the number of steps in output $i$. These step-level rewards are then standardized by their global mean and standard deviation: $\widetilde{r}_i^{index(j)} = \frac{r_i^{index(j)} - {\rm mean(\mathbf{R})}}{{\rm std(\mathbf{R})}}$. The process supervision method determines the token-wise advantage $\hat{A}_{i, t}$ as the aggregate of normalized rewards for all subsequent steps, i.e., $\hat{A}_{i, t} = \sum_{index(j) \ge

\subsubsection{Iterative RL with GRPO} During reinforcement learning, the initially trained reward model may become inadequate for effectively guiding the evolving policy model. To address this, we investigate an iterative approach to RL utilizing GRPO. As detailed in Algorithm \ref{alg:iter-grpo}, our iterative GRPO framework constructs new datasets for the reward model by leveraging samples generated from the current policy model. The reward model is continuously refined through a replay mechanism that integrates 10\% of prior data into the training batches. Subsequently, the updated policy model becomes the new reference model, and training of the policy model proceeds using the newly updated reward model.

\subsection{Training and Evaluating DeepSeekMath-RL}

Our RL experiments use DeepSeekMath-Instruct 7B as the base model. The reinforcement learning data are chain-of-thought questions drawn from GSM8K and MATH subsets within the SFT dataset, amounting to approximately 144K instances. We intentionally omit other SFT-derived questions to better analyze RL’s effects on benchmarks where training data is absent throughout RL. Construction of reward model training data follows the protocol from \citep{wang2023math}. The initial reward model is trained on DeepSeekMath-Base 7B, employing a learning rate of 2e-5. In GRPO, the policy model uses a learning rate of 1e-6 and a KL penalty of 0.04. For each training sample, we generate $64$ completions. Sequences are constrained to a maximum length of 1024 tokens, and batches comprise 1024 samples. Following each exploration round, the policy model is updated once. Evaluation protocols for DeepSeekMath-RL 7B follow those established for DeepSeekMath-Instruct 7B. Here, GSM8K and MATH with chain-of-thought annotations are categorized as in-domain tasks, while all remaining benchmarks are considered out-of-domain.

Table \ref{tab:sft_rl_math} presents the comparative results of open- and closed-source models employing chain-of-thought as well as tool-integrated reasoning approaches on English and Chinese tasks. Our observations are as follows: 1) DeepSeekMath-RL 7B achieves 88.2\% accuracy on GSM8K and 51.7\% accuracy on MATH through chain-of-thought reasoning, outperforming every open-source model between 7B and 70B in size, and also exceeding the results of most proprietary models. 2) Importantly, DeepSeekMath-RL 7B is exclusively trained on chain-of-thought style instruction tuning data drawn from GSM8K and MATH, and is initialized from DeepSeekMath-Instruct 7B. Although its training data are highly specialized, it demonstrates improvements over DeepSeekMath-Instruct 7B in all tested aspects, underscoring the impact of reinforcement learning.

\section{Discussion}
Here, we discuss insights derived from both our pre-training and RL experimental processes.

\subsection{Lessons Learnt in Pre-Training}

We begin by elaborating on our experiences with pre-training. Except where noted, we follow the experimental parameters described in Section \ref{sec:quality-policy}. For clarity, references to the DeepSeekMath Corpus in this part denote the version containing 89B tokens produced during the second data collection cycle.

\subsubsection{Code Training Benefits Mathematical Reasoning}

A commonly held yet unconfirmed belief posits that exposure to code during training enhances reasoning capabilities. We seek to partially address this, with a particular emphasis on mathematical reasoning: training with code data strengthens a model’s performance in mathematical reasoning tasks, whether or not external tools are applied.

To examine the impact of code-oriented training on mathematical reasoning ability, we conducted both two-stage and one-stage training experiments as detailed below:

\textbf{Two-Stage Training}

\begin{itemize}[topsep=0pt]
    \item \textbf{Code Training for 400B Tokens $\rightarrow$ Math Training for 150B Tokens}: DeepSeek-LLM 1.3B is pre-trained on 400B code tokens, followed by a further 150B tokens of mathematical data;
    \item \textbf{General Training for 400B Tokens $\rightarrow$ Math Training for 150B Tokens}: For comparison, we replace code tokens with general domain tokens (sampled from a broad and diverse corpus curated by DeepSeek-AI) during the initial phase, allowing us to assess the distinct benefits of code-centric pretraining versus general textual data for boosting mathematical reasoning ability.
\end{itemize}

\textbf{One-Stage Training}

\begin{itemize}[topsep=0pt]
    \item \textbf{Math Training for 150B Tokens}: Here, DeepSeek-LLM 1.3B is exclusively trained on 150B math tokens;
    \item \textbf{Training on a mixture of 400B Code Tokens and 150B Math Tokens}: As code-first, math-second training leads to diminished coding skill, we further assess a training strategy where code and math tokens are interleaved from the outset, aiming to determine whether such a mix simultaneously enhances mathematical reasoning and counters catastrophic forgetting.
\end{itemize}

\paragraph{Results}
Empirical findings presented in Table \ref{tab:code-math} and Table \ref{tab:code-math-general-eval} report how downstream performance shifts across the evaluated training regimes.

Pretraining on code data proves advantageous for program-based mathematical reasoning in both multi-stage and single-stage training contexts. Specifically, Table \ref{tab:code-math} indicates that in the two-phase scenario, code-only pretraining yields marked gains in solving tasks from GSM8K and MATH datasets using Python, with further increases upon subsequent math-focused fine-tuning. Notably, when code and math tokens are trained together in a single-stage process, this integration alleviates the catastrophic forgetting typically observed in staged setups, and produces additive improvements in both coding proficiency (Table \ref{tab:code-math-general-eval}) and in program-aided mathematical reasoning (Table \ref{tab:code-math}).

Beyond program-facilitated reasoning, code-based pretraining also enhances math problem solving without tools. In two-stage training, an initial phase of code exposure offers moderate but tangible benefits and fosters better adaptation during follow-up math learning, eventually attaining optimal performance. By contrast, mixing code and math tokens from the start in single-stage training appears to hinder non-tool-based mathematical reasoning. This outcome may be explained by the restricted model capacity of DeepSeek-LLM 1.3B, which may be insufficient to jointly master both code and mathematics domains when trained on them simultaneously.

\subsubsection{ArXiv Papers Appear to Offer Limited Benefits for Enhancing Mathematical Reasoning}

ArXiv papers frequently form a substantial portion of math pre-training datasets \citep{minerva,gpt-f,llemma,mathpile}, yet their concrete effects on enhancing mathematical reasoning remain largely under-examined. Contrary to prevalent assumptions, our experimental findings suggest that arXiv papers do not substantially contribute to the advancement of mathematical reasoning skills. We evaluated this by training models of various scales, specifically DeepSeek-LLM 1.3B and DeepSeek-Coder-Base-v1.5 7B \citep{deepseek-coder}, utilizing two distinct arXiv datasets processed with different pipelines:

\begin{itemize}[topsep=0pt]
    \item \textbf{MathPile} \citep{mathpile}: An 8.9B-token dataset curated with rigorous cleaning and filtering heuristics, composed of over 85\% scientific content sourced from arXiv;
    \item \textbf{ArXiv-RedPajama} \citep{redpajama}: Contains the full set of arXiv LaTeX files after removal of preambles, comments, macros, and bibliographies, yielding a total of 28.0B tokens.
\end{itemize}

For our experiments, we separately trained DeepSeek-LLM 1.3B for 150B tokens and DeepSeek-Coder-Base-v1.5 7B for 40B tokens on each of these arXiv-based datasets. The results consistently indicate that exposure solely to arXiv papers does not translate into observable improvements— and in some cases leads to a decrease—in performance across several mathematical reasoning benchmarks spanning multiple complexity levels. This evaluation utilized quantitative problem-solving datasets such as GSM8K and MATH (see Table \ref{tab:arxiv-cot}), multiple-choice assessment tasks like MMLU-STEM (Table \ref{tab:arxiv-cot}), as well as formal mathematics datasets including miniF2F (Table \ref{tab:arxiv_atp}).

Nonetheless, it is important to recognize that our findings are subject to certain caveats. Specifically, we did not assess:

\begin{itemize}[topsep=0pt]
    \item The potential utility of arXiv-derived data for specialized math-related objectives absent from our benchmark suite, e.g., translating formalized mathematical arguments into informal language;
    \item The interaction effects when arXiv data is mixed with alternative sources of pre-training data;
    \item The possibility that benefits from arXiv papers could arise when training larger-scale models.
\end{itemize}

Accordingly, a more comprehensive investigation is warranted to clarify these open questions, which we leave as directions for subsequent research.

\subsection{Perspectives on Reinforcement Learning}
\subsubsection{Moving Toward a Unified Framework} Here, we outline a unifying perspective for analyzing diverse training techniques—namely SFT, RFT, DPO, PPO, and GRPO—and conduct further studies to dissect key elements within this general framework. The gradient of a given training algorithm with respect to its parameter $\theta$ can in general be formalized as:

\begin{equation}
    \nabla_{\theta}\mathcal{J}_{\textcolor{red}{\mathcal{A}}}(\theta) = \mathbb{E}[\underbrace{(q,o) \sim \textcolor{red}{\mathcal{D}}}_{Data \ Source}]\left( \frac{1}{|o|} \sum_{t=1}^{|o|}  \underbrace{GC_{{\mathcal{A}}}(q, o, t, \textcolor{red}{\pi_{{rf}}})}_{Gradient \ Coefficient}  \nabla_{\theta}\log \pi_{\theta}(o_t | q, o_{<t})\right).
\label{eq:objective}
\end{equation}

This formulation highlights three principal elements: (1) the \textit{Data Source $\mathcal{D}$}, which specifies the origin of the data used for training; (2) the \textit{Reward Function $\pi_{{rf}}$}, supplying the reward signal that guides model updates; and (3) the \textit{Algorithm $\mathcal{A}$}, responsible for transforming both training data and the reward signal into the gradient coefficient $GC$, which in turn influences the strength and direction of learning (i.e., how strongly a sample is penalized or reinforced). Using this unified perspective, we discuss the characteristics of several widely used training strategies:

\begin{itemize}
    \item  \textbf{Supervised Fine-tuning (SFT)}: In SFT, a pretrained model undergoes additional training using datasets curated and labeled by humans.
    \item \textbf{Rejection Sampling Fine-tuning (RFT)}: RFT builds upon SFT by retraining the SFT model, focusing on its outputs that have been filtered according to correctness, which are sampled in response to SFT prompts.
    \item \textbf{Direct Preference Optimization (DPO)}: DPO further tunes the SFT model by leveraging a loss function based on pairwise preferences, utilizing an augmented set of outputs generated from the SFT model.
    \item \textbf{Online Rejection Sampling Fine-tuning (Online RFT)}: In contrast to RFT, Online RFT initializes its policy model from SFT and incrementally fine-tunes it with new samples produced by the continually updated policy model itself.
    \item \textbf{PPO/GRPO}: In these approaches, the policy model is seeded with the SFT model and is progressively enhanced by reinforcement learning, using outputs directly sampled from the current version of the policy model.
\end{itemize}

The constituent elements of these methodologies are outlined in Table \ref{tab:data_policy}. For an in-depth exposition of the derivation steps, consult Appendix \ref{app:analysis-rl}.

\paragraph{Analysis of Data Source} We classify data sources into two principal groups: online sampling and offline sampling. Online sampling refers to instances where the training set is drawn from outputs generated by the contemporaneous training policy model, whereas offline sampling utilizes data produced by the original SFT model. RFT and DPO are representative of offline methods, whereas Online RFT and GRPO are characterized by online sampling.

Referencing Figure \ref{fig:rl-analysis}, our analysis reveals that Online RFT substantially outperforms RFT across both benchmarks assessed. During the early phases of training, performance is largely equivalent between Online RFT and RFT; however, Online RFT establishes a clear superiority as training proceeds, underscoring the efficacy of online approaches. This observation can be rationalized as follows: initially, the policy model (actor) and the SFT model are highly similar, which results in training data with minimal variation. As training advances, however, the divergence between the actor's output and the initial model increases, enabling online data sampling to yield substantially improved results due to increased diversity and relevance.

\paragraph{Analysis of Gradient Coefficient} The mechanism by which the reward signal is translated into a gradient coefficient for model parameter updates is critical. In our experiments, the reward function is operationalized in two forms: `Rule' and `Model.' The Rule-based approach assesses response quality based on the accuracy of the answer, while the Model-based approach employs a trained reward model—whose training data is annotated via the rule method—to assign scores to each response. A central distinction, as evidenced by Equations \ref{eq:GC-RFT} and \ref{eq:GC-GRPO}, is that GRPO dynamically modulates its gradient coefficient according to the specific reward score from the reward model. This confers the advantage of rewarding or penalizing responses in proportion to their merit. Online RFT, in contrast, does not implement such penalization; it uniformly reinforces all correct responses, regardless of any variance in response quality.

As illustrated in Figure \ref{fig:rl-analysis}, GRPO consistently outperforms online RFT, underlining the effectiveness of manipulating the positive and negative gradient coefficients. Moreover, GRPO+PS achieves better results than GRPO+OS, demonstrating the advantages of employing more granular, step-sensitive gradient coefficients. We additionally investigate iterative RL in our study by executing two iterations. Figure \ref{fig:iter-rl} reveals that iterative RL provides notable performance gains, with the most substantial improvement occurring after the initial iteration.

\subsubsection{Why RL Works?} This study implements reinforcement learning using a selection of data from instruction tuning, which leads to substantial improvements relative to the instruction-tuned model. To gain a deeper understanding of why RL is beneficial, we assess Pass@K and Maj@K accuracy for both the Instruct and RL models across two benchmarks. Figure \ref{fig:maj-pass} indicates that RL leads to better Maj@K outcomes, though it does not affect Pass@K. These observations suggest that RL strengthens overall model performance by increasing the robustness of the output distribution—in essence, \textbf{the observed gains appear to arise from elevating the likelihood of correct responses among the TopK options, rather than from a boost in core abilities.} In a similar vein, \citep{wang2023making} describes a \textbf{misalignment problem} encountered in reasoning tasks using the SFT model and demonstrates that preference alignment strategies can successfully enhance the reasoning abilities of SFT models \citep{yuan2023rrhf,song2023preference,wang2023making}.

\subsubsection{How to Achieve More Effective RL?}
\label{sec:effec_rl}
Our findings confirm that RL is highly effective for mathematical reasoning tasks. Additionally, we introduce a unified framework that contextualizes various mainstream training methodologies, positing that they can be viewed as either direct forms of RL or simplified derivatives. According to Equation \ref{eq:objective}, the paradigm consists of three principal elements: Data Source, Algorithm, and Reward Function. We propose future research opportunities for each component.

\paragraph{Data Source} The data source serves as the foundational input for all training frameworks. In the RL context, it specifically denotes unlabeled queries paired with responses sampled from the policy model. Our approach in this work utilizes only those questions seen during instruction tuning and applies straightforward nucleus sampling for output generation. We posit that this approach might explain why our RL pipeline exclusively enhances Maj@K accuracy. For future work, we aim to extend our RL pipeline to prompts drawn from distributions beyond the training data, while experimenting with \textbf{advanced sampling (decoding) strategies} such as tree-search-based approaches \citep{yao2023tree}. Furthermore, \textbf{efficient inference techniques} \citep{xia-etal-2023-speculative,leviathan2023fast,kwon2023efficient,xia2024unlocking}, which greatly influence the exploration capability of policy models, remain a critical area for advancement.

\paragraph{Algorithms} Algorithms utilize both data and reward signals to compute the gradient coefficients necessary for updating model parameters. According to Equation \ref{eq:objective}, contemporary approaches largely \textbf{TRUST} the reward function to modulate the conditional probability of generating specific tokens. Nonetheless, the reward signal's accuracy cannot be guaranteed at all times, particularly when tackling highly intricate tasks. For instance, the PRM800K datasets \citep{lightman2023let}, despite meticulous annotation by expert annotators, still contain roughly 20\% mislabeled entries\footnote{\url{https://github.com/openai/prm800k/issues/12\#issuecomment-1728491852}}. Consequently, there is a pressing need to develop reinforcement learning algorithms that maintain robustness in the presence of noisy reward signals. We propose that alignment strategies following the \textbf{WEAK-TO-STRONG} paradigm \citep{burns2023weak} have the potential to fundamentally transform learning algorithm frameworks.

\paragraph{Reward Function} The reward function provides the core feedback used for training. In the context of reinforcement learning, this function is typically implemented as a neural reward model. We identify three primary research trajectories for advancing reward models: 1) \textbf{Improving the generalization capacity of reward models.} Effective reward models should generalize beyond their training distributions, accommodating novel queries and advanced model outputs; if they do not, reinforcement learning risks merely reinforcing the status quo rather than enhancing underlying LLM performance; 2) \textbf{Incorporating and quantifying reward model uncertainty.} Characterizing uncertainty can serve as an interface connecting imperfect reward models to weak-to-strong learning paradigms; 3) \textbf{Constructing efficient, high-caliber process reward models} that are able to deliver detailed, step-wise feedback throughout reasoning tasks \citep{lightman2023let,wang2023math}.

\section{Conclusion, Limitation, and Future Work}

In this work, we introduce DeepSeekMath, which establishes state-of-the-art results among open-source models on the MATH benchmark designed for competition-level tasks, and narrows the performance gap with proprietary models. The model is built upon DeepSeek-Coder-v1.5 7B and is further continually trained with 500B tokens—of which 120B are mathematical tokens predominantly drawn from Common Crawl. Our comprehensive ablation experiments reveal that web-based resources are particularly promising as a source of rich mathematical content, whereas the anticipated advantages of arXiv appear to be limited. We put forth Group Relative Policy Optimization (GRPO), extending the conventional Proximal Policy Optimization (PPO), to significantly strengthen mathematical reasoning while also reducing memory requirements. Our results confirm that GRPO continues to be beneficial even for DeepSeekMath-Instruct 7B at high benchmark performance levels. Moreover, we propose a consolidated perspective on related methodologies and outline several research avenues to further improve reinforcement learning strategies.

Despite these advancements, DeepSeekMath~shows certain shortcomings, especially in addressing geometry and theorem-proving tasks when compared to closed-source models. For example, during preliminary assessments, the model was unable to resolve problems related to triangles and ellipses—suggesting possible biases introduced during the data curation processes in pre-training and fine-tuning. Additionally, the model’s performance in few-shot scenarios lags behind that of GPT-4, which demonstrates marked improvement with few-shot prompts, whereas DeepSeekMath~exhibits almost no difference between zero-shot and few-shot results. Moving forward, we plan to refine our data curation pipeline to assemble an even higher-quality pre-training corpus. Furthermore, we will investigate the promising directions discussed in Section \ref{sec:effec_rl} to enhance reinforcement learning methodologies for large language models.

\end{document}